% Standalone mathematical sections for inclusion in the common subcritical paper.
% Authoring agent: /root/subcritical_sample_gp, 2026-10-15.
% Requires amsmath, amssymb, amsthm; uses theorem, lemma, proposition environments.

\section{The Gaussian sign posterior: \texorpdfstring{$\gamma<0$}{gamma < 0}}
\label{sec:negative-gamma}

Write $M=2^d$, let $K$ denote the $M\times M$ limiting covariance
matrix, and let $v=M\lambda_d>0$ be its smallest matrix eigenvalue.
All norms and averages on the cube below use the uniform probability
measure.  The RKHS norm of a vector $h$ is
$\|h\|_K^2=h^TK^{-1}h$.
Let $S$ be any training sample, let $s$ be its number of distinct
locations, and condition $G\sim N(0,K)$ on the observed sign constraints.
In the regime $\gamma<0$ this is exactly the limiting posterior: the
fixed original margin vanishes after output normalization.

\begin{lemma}[A norm bound for the conditional mean]
\label{lem:sign-mean-norm}
Let $m(x)=\mathbb E[G(x)\mid S]$.  Then
\[
 m=K_{\mathcal X,S}K_{S,S}^{-1}m_S,
 \qquad \|m\|_K^2\le s.
\]
\end{lemma}
\begin{proof}
Distinct training coordinates have a nonsingular covariance matrix.
Whiten their Gaussian vector as $G_S=K_{S,S}^{1/2}\xi$, with
$\xi\sim N(0,I_s)$.  The observed sign event is a homogeneous cone
in $\xi$.  A standard Gaussian radius is independent of its direction,
so conditioning on this cone leaves
$\mathbb E[\|\xi\|^2\mid S]=s$.  Jensen's inequality gives
$m_S^TK_{S,S}^{-1}m_S\le s$.  Gaussian conditioning gives the
representer formula; its RKHS norm squared is exactly the preceding
quadratic form.
\end{proof}

\begin{lemma}[Posterior confidence gives a latent margin]
\label{lem:confidence-latent-margin}
Put
\[
 \theta_0=\frac{13}{20},\qquad
 c_* = \sqrt{2\pi}\left(\frac1e-\frac7{20}\right)>0.
\]
At any unobserved input $x$, correct-label probability strictly
larger than $\theta_0$ implies
\[
 f(x)m(x)>c_*\sqrt v.
\]
\end{lemma}
\begin{proof}
The conditioned Gaussian is log-concave, since its conditioning
set is convex.  Its scalar marginal $X=f(x)G(x)$ is therefore
log-concave.  If $\mu=\mathbb E X$, both its distribution function
$F_X$ and survival function are log-concave.  Jensen's inequality
and the probability integral transform give
\[
 F_X(\mu)\ge e^{-1},\qquad
 \mathbb P(X\ge\mu)\ge e^{-1}.
\]
For example,
$\log F_X(\mu)\ge\mathbb E\log F_X(X)=-1$.
Consequently $\mathbb P(X>0)>\theta_0$ implies $\mu>0$,
since $\theta_0>1-e^{-1}$.

Conditionally on the full vector $G_S$, $G(x)$ is Gaussian with
variance
\[
 \sigma_x^2=K(x,x)-K(x,S)K_{S,S}^{-1}K(S,x).
\]
The same variance occurs in every component after conditioning only
on the signs, so the scalar density is at most
$(\sqrt{2\pi}\sigma_x)^{-1}$.  Since
$\mathbb P(X>\mu)\le1-e^{-1}$, we obtain
\[
 \theta_0<\mathbb P(X>0)
 \le1-e^{-1}+\frac{\mu}{\sqrt{2\pi}\sigma_x},
\]
which implies $\mu>c_*\sigma_x$.
Finally, $K\succeq vI$, and for every vector $b$ of training
coefficients,
\[
 \operatorname{Var}(G(x)-b^TG_S)
 \ge v(1+\|b\|_2^2)\ge v.
\]
Minimization over $b$ gives $\sigma_x^2\ge v$.
\end{proof}

Define the polynomial-threshold parity obstruction
\[
 B(d,r)=2^{-d}\sum_{j=0}^{\lfloor(d-r-1)/2\rfloor}\binom dj
 \quad(0\le r<d).
\]
The polynomial-threshold lemma proved above says that every real
polynomial of degree at most $r$, thresholded at zero, has error
at least $B(d,r)$ against full parity.  The convention at zero can
be either fixed sign: adding an arbitrarily small constant removes
all zeros without changing that convention.

\begin{theorem}[A confidence bound for negative exponents]
\label{thm:negative-confidence}
For full parity, every training sample with $s$ distinct locations
satisfies, in the $\gamma<0$ width limit,
\begin{equation}\label{eq:negative-confidence}
 F_{\theta_0}(S)\ge B(d,r)-\frac{s}{M}
             -\frac{s\lambda_{r+1}}{c_*^2v},
 \qquad 0\le r<d.
\end{equation}
Here $F_{\theta_0}(S)$ is the fraction of inputs where correct-label
probability is at most $\theta_0=13/20$.  These points also fail
the requested strict $2/3$ confidence criterion.  Moreover, for the
original finite-width posteriors,
\[
 \liminf_{N\to\infty}F_N(S)
 \ge B(d,r)-\frac{s}{M}
             -\frac{s\lambda_{r+1}}{c_*^2v}.
\]
\end{theorem}
\begin{proof}
Let $p$ be the Fourier truncation of $m$ to degree at most $r$.
The eigenvalues decrease with Fourier degree, and
Lemma~\ref{lem:sign-mean-norm} gives
\[
 \|m-p\|_2^2\le\lambda_{r+1}\|m\|_K^2
                    \le s\lambda_{r+1}.
\]
At an unobserved input with confidence strictly larger than $\theta_0$,
Lemma~\ref{lem:confidence-latent-margin} gives
$f(x)m(x)>c_*\sqrt v$.  If the sign of $p(x)$ is wrong there,
then $|m(x)-p(x)|\ge c_*\sqrt v$.  Thus the error of the
polynomial threshold $p$ is at most
$F_{\theta_0}(S)+s/M+s\lambda_{r+1}/(c_*^2v)$.
The polynomial-threshold lower bound proves the claim.  The term
$s/M$ is needed because observed coordinates do not have a
positive Gaussian residual variance.
\end{proof}

\section{A fixed global margin: \texorpdfstring{$\gamma=0$}{gamma = 0}}
\label{sec:zero-gamma}

Now the limiting posterior is $G\sim N(0,K)$ conditioned on
$|G(x)|>a$ at every input and on the observed labels.
This conditioning is not convex.  The argument in the previous
section therefore does not apply.  Instead we control the Fourier
tail of an entire posterior draw.

\begin{lemma}[High-degree covariance mass]
\label{lem:tail-trace}
Put $\beta=2^{-(L-1)}$ and let
$\tau_r=\sum_{j>r}\binom dj\lambda_j$.  For every $r\ge0$,
\[
 \tau_r\le \frac{2}{1-e^{-1}}(r+1)^{-\beta}.
\]
\end{lemma}
\begin{proof}
Write the saturated kernel function as
$F(t)=H^{\circ(L-1)}(t)=\sum_{m\ge0}c_mt^m$, where $c_m\ge0$
and $\sum_m c_m=1$.  The kernel
$((1+x\cdot y)/(d+1))^m$ has Fourier degree at most $m$ and
normalized trace one.  Consequently
$\tau_r\le\sum_{m>r}c_m$.
For $u=1/(r+1)$,
\[
 (1-e^{-1})\sum_{m>r}c_m
 \le1-F(1-u).
\]
The elementary inequality
$1-H(1-u)\le\sqrt{2u}$, valid for $0\le u\le1$, follows
from $u=2\sin^2(\tfrac12\arccos(1-u))$ and
$\sin t\ge2t/\pi$ on $[0,\pi/2]$.
Iteration gives $1-F(1-u)\le2u^{\beta}$.
\end{proof}

\begin{lemma}[A denominator bound uniform in the labels]
\label{lem:margin-denominator}
Set
\[
 t=\sqrt{2\log(4M)},\qquad b=a+t+1,\qquad
 D_s=b^2\left(\lambda_0^{-1}+\frac{4s}{v}\right).
\]
The event $A_S$ imposing the global margin and any consistent labels
on $s$ distinct inputs satisfies
\[
 \mathbb P(A_S)\ge\tfrac12 e^{-D_s}.
\]
\end{lemma}
\begin{proof}
Since every marginal of $G$ has variance one, a union bound gives
$\mathbb P(E)\ge1/2$ for the symmetric set
$E=\{\max_x|G(x)|\le t\}$.
Let $u$ be the indicator vector of the observed locations with
negative labels, and put $h=b(\mathbf1-2u)$.  Then $h+E\subset A_S$.
The constant vector has squared RKHS norm $1/\lambda_0$, while
$\|u\|_K^2\le s/v$.  Thus
\[
 \tfrac12\|h\|_K^2
 \le b^2(\lambda_0^{-1}+4s/v)=D_s.
\]
The Gaussian shift identity and symmetry of $E$ give
\begin{align*}
 \mathbb P(G\in h+E)
 &=e^{-\|h\|_K^2/2}
   \mathbb E[\mathbf1_E(G)e^{-h^TK^{-1}G}]\\
 &\ge e^{-\|h\|_K^2/2}\mathbb P(E).
\end{align*}
Indeed the conditional mean of the linear form on the symmetric
set $E$ is zero, so Jensen's inequality applies.  This proves
the bound.
\end{proof}

\begin{theorem}[A confidence bound at zero exponent]
\label{thm:zero-confidence}
Fix $0\le r<d$ and $\delta>0$ with $\tau_r\le\delta/4$.
Put $\ell=\lambda_{r+1}$ and
\[
 q_s=\min\{1,\ 2\exp(D_s-\delta/(8\ell))\}.
\]
For full parity, the limiting posterior mean error $R(S)$ and
confidence failure $F(S)$ satisfy
\begin{align}
 R(S)&\ge B(d,r)-\delta/a^2-q_s,\label{eq:zero-risk}\\
 F(S)&\ge\max\left\{0,
       \frac{3[B(d,r)-\delta/a^2-q_s]-1}{2}\right\}.
       \label{eq:zero-confidence}
\end{align}
\end{theorem}
\begin{proof}
Let $G_{>r}$ be the high-degree Fourier part of the Gaussian draw.
In an orthonormal Fourier basis,
$W=\|G_{>r}\|_2^2=\sum_{|U|>r}\lambda_{|U|}\xi_U^2$,
with independent standard Gaussian $\xi_U$.
Every coefficient is at most $\ell$, and their sum is $\tau_r$.
For $\theta=1/(4\ell)$,
\[
 \log\mathbb E e^{\theta W}
 =-\tfrac12\sum_{|U|>r}\log(1-2\theta\lambda_{|U|})
 \le2\theta\tau_r.
\]
Markov's inequality and $\tau_r\le\delta/4$ imply
\[
 \mathbb P(W>\delta)\le e^{-\delta/(8\ell)}.
\]
Lemma~\ref{lem:margin-denominator} therefore gives
$\mathbb P(W>\delta\mid A_S)\le q_s$.

On $A_S$, if $\operatorname{sign}G$ and the threshold of
$G_{\le r}$ disagree at an input, then $|G_{>r}(x)|>a$.
When $W\le\delta$, their disagreement fraction is consequently
at most $\delta/a^2$.  The degree-$r$ polynomial threshold has
parity error at least $B(d,r)$.  Hence the posterior draw has
error at least $B(d,r)-\delta/a^2$ except on an event of
probability at most $q_s$.  Averaging gives
\eqref{eq:zero-risk}.
Finally, if $p(x)$ is correct-label posterior probability, then
$R=\mathbb E_x(1-p(x))\le (1-F)/3+F$.
Rearranging proves~\eqref{eq:zero-confidence}.
\end{proof}

\subsection{Consequence at fixed depth}
For fixed $L$ and fixed margin $a>0$, the preceding bounds imply
constants $c>0$ and $d_0$ such that, for $d\ge d_0$, every sample
with
\[
 n\le\exp(c\sqrt d\log(d+1))
\]
has limiting confidence failure at least $1/16$, for every
$\gamma\le0$.
To see this, take $r=\lfloor\sqrt d/64\rfloor$.
Elementary central-binomial estimates give $B(d,r)\ge7/16$ for
all sufficiently large $d$.  Spectral ordering and trace one give
$\lambda_{r+1}\le\binom d{r+1}^{-1}
=\exp[-\Omega(\sqrt d\log(d+1))]$.
The spectral bounds proved earlier give
$\lambda_0^{-1}=O_L(d)$ and $v^{-1}=O_L(d\log(d+1))$.
For negative exponents these estimates apply directly to
\eqref{eq:negative-confidence}; its two error terms tend to zero
when $c$ is sufficiently small.
At exponent zero choose $\delta=a^2/32$.
Lemma~\ref{lem:tail-trace} eventually gives $\tau_r\le\delta/4$,
and the same estimates ensure $D_n\le\delta/(16\ell)$ and
$q_n\le1/32$.  Thus $R\ge3/8$ and $F\ge1/16$.
The threshold $d_0$ can be very large at large fixed depth because
the trace exponent is $2^{-(L-1)}$.

These are bounds for the limiting posterior, and they also give
$\liminf_{N\to\infty}F_N(S)\ge1/16$ for the requested finite-width
strict $2/3$ criterion.  At negative exponents, the proof uses
$\theta_0=13/20<2/3$: every input counted by limiting
$F_{\theta_0}$ remains a failure at all sufficiently large widths.
The input cube is finite, so this holds simultaneously on the
counted inputs.  At exponent zero, the risk lower bound is continuous
under finite-state posterior convergence; the universal inequality
$F_N\ge(3R_N-1)/2$ transfers it directly.
